% Chapter 1 worked example. Included by ch01_resources.tex.
\section{Worked Example: Emergency-Triage Intervention Portfolio}
\label{sec:comp-ch01-worked-example}

\subsection{Purpose and Status}
This worked example demonstrates how the methods in Chapter~1 can be
translated into a reviewable institutional decision. All numerical inputs are
synthetic unless a source is identified explicitly. They are intended to show the
logic of the analysis, not to report empirical findings or establish a universal
threshold. Local use requires replacement of the assumptions with current,
versioned evidence and approval through the relevant clinical, managerial,
economic, legal, technical, and governance processes.

A regional public hospital is considering an AI-supported emergency-triage service after repeated crowding, prolonged time to first clinical assessment, and concern about delayed recognition of deteriorating patients. The proposal must be assessed as a complete intervention rather than as a prediction model. 

\subsection{Decision Problem and Comparators}
\textbf{Decision problem.} Should the hospital pilot an AI-supported triage pathway, strengthen the current pathway through staffing and protocol redesign, or adopt a lower-complexity rules-based alternative?

The comparator set is deliberately broader than the existing pathway. This prevents
a new technology from receiving a lower evidentiary threshold than staffing,
workflow, shared-service, conventional digital, or no-change alternatives.

\begin{longtable}{@{}p{0.08\textwidth}p{0.84\textwidth}@{}}
\caption{Comparators for the Chapter 1 worked example}
\label{tab:comp-ch01-comparators}\\
\toprule
\textbf{Option} & \textbf{Comparator or strategy} \\
\midrule
\endfirsthead
\toprule
\textbf{Option} & \textbf{Comparator or strategy} \\
\midrule
\endhead
\bottomrule
\endlastfoot
1 & Current triage with existing staffing \\
2 & Additional triage staffing during peak periods \\
3 & Registration and streaming redesign \\
4 & A conventional rules-based early-warning score \\
5 & AI-supported prioritisation with human confirmation \\
\end{longtable}

\subsection{Model and Decision Structure}
The analytical structure should follow the decision rather than the available
software. The team should map the causal pathway from intervention input to
information, human decision, action, outcome, resource consequence, and review.
The structure should also identify capacity constraints, uncertainty, subgroup
variation, implementation requirements, and events that would change the
recommendation.

A compact quantitative relationship used in this example is:
\begin{equation}
\text{Incremental value hypothesis: }\quad \Delta V = f(\Delta O,\,\Delta C,\,\Delta W,\,\Delta R)
\label{eq:comp-ch01-key-relationship}
\end{equation}
The equation is an organizing aid rather than a substitute for disaggregated evidence,
qualitative findings, or mandatory safeguards. Its variables and interpretation must
be defined in the local protocol.

\subsection{Synthetic Parameters and Evidence Sources}
Table~\ref{tab:comp-ch01-parameters} provides the base-case inputs. A
production analysis should add parameter distributions, correlations, structural
alternatives, data dates, system versions, and evidence-confidence ratings.

\begin{longtable}{@{}p{0.32\textwidth}p{0.22\textwidth}p{0.38\textwidth}@{}}
\caption{Synthetic parameters for the Chapter 1 worked example}
\label{tab:comp-ch01-parameters}\\
\toprule
\textbf{Parameter} & \textbf{Base value} & \textbf{Source or interpretation} \\
\midrule
\endfirsthead
\toprule
\textbf{Parameter} & \textbf{Base value} & \textbf{Source or interpretation} \\
\midrule
\endhead
\bottomrule
\endlastfoot
Annual eligible attendances & 72,000 & Local activity report \\
Baseline median time to first assessment & 84 minutes & Emergency-department dashboard \\
High-risk event prevalence & 4.5\% & Retrospective cohort \\
AI sensitivity at proposed threshold & 0.88 & External evidence; local verification required \\
AI specificity at proposed threshold & 0.72 & External evidence; local verification required \\
Incremental annual programme cost & EUR 310,000 & Local lifecycle-cost estimate \\
Available additional rapid-assessment capacity & 18 patients/day & Workforce and space assessment \\
\end{longtable}

\subsection{Step-by-Step Analysis}
The sequence below is intended to make the analysis reproducible and to ensure
that evidence, economics, governance, and implementation develop together.

\begin{longtable}{@{}p{0.07\textwidth}p{0.55\textwidth}p{0.30\textwidth}@{}}
\caption{Analytical sequence}
\label{tab:comp-ch01-analysis-steps}\\
\toprule
\textbf{Step} & \textbf{Required work} & \textbf{Decision record} \\
\midrule
\endfirsthead
\toprule
\textbf{Step} & \textbf{Required work} & \textbf{Decision record} \\
\midrule
\endhead
\bottomrule
\endlastfoot
1 & Map the present pathway and identify the decision window in which prioritisation could still change care. & Evidence and responsible owner to be recorded \\
2 & Compare all realistic alternatives before assigning value to the AI option. & Evidence and responsible owner to be recorded \\
3 & Represent the intervention bundle: data availability, interface, triage-nurse review, escalation, rapid-assessment capacity, monitoring, and fallback. & Evidence and responsible owner to be recorded \\
4 & Estimate false-positive workload and false-negative consequences at the proposed threshold. & Evidence and responsible owner to be recorded \\
5 & Test whether downstream capacity can absorb the additional patients identified as high risk. & Evidence and responsible owner to be recorded \\
6 & Define a reversible pilot with clinical, operational, economic, equity, and trust indicators. & Evidence and responsible owner to be recorded \\
\end{longtable}

\subsection{Illustrative Outputs}
The outputs in Table~\ref{tab:comp-ch01-outputs} show how technical,
clinical, operational, economic, equity, and governance findings should be kept
visible rather than compressed into a single favourable or unfavourable score.

\begin{longtable}{@{}p{0.30\textwidth}p{0.24\textwidth}p{0.38\textwidth}@{}}
\caption{Illustrative model and decision outputs}
\label{tab:comp-ch01-outputs}\\
\toprule
\textbf{Output} & \textbf{Result} & \textbf{Interpretation} \\
\midrule
\endfirsthead
\toprule
\textbf{Output} & \textbf{Result} & \textbf{Interpretation} \\
\midrule
\endhead
\bottomrule
\endlastfoot
Expected high-risk alerts per day & Approximately 67 & Includes true and false positive alerts \\
Additional patients requiring rapid review & Approximately 35/day & Exceeds currently available capacity without redesign \\
Primary value mechanism & Earlier review of high-risk patients & Dependent on response capacity \\
Main economic risk & Unusable alerts and displaced staff time & May increase rather than reduce delay \\
Preferred initial decision & Redesign and limited pilot & Not unrestricted deployment \\
\end{longtable}

\subsection{Sensitivity, Uncertainty, and Subgroups}
At minimum, the completed analysis should vary the parameters most likely to
change the decision, test at least one credible alternative model structure, and
report subgroup results separately where performance, access, response, burden,
or opportunity cost may differ. Deterministic analysis should identify thresholds;
probabilistic analysis should be used where credible parameter distributions and
correlations are available. Scenario analysis is preferable when structural or
future uncertainty cannot be represented defensibly by one probability model.

The record should distinguish uncertainty that can be reduced through evidence,
uncertainty that can be managed through safeguards, and uncertainty that remains
too consequential to accept. Missing evidence should not be represented as an
average or neutral value without justification.

\subsection{Interpretation and Recommendation}
Proceed only with a redesigned, time-limited pilot that adds response capacity and compares the AI pathway with a staffing-and-streaming alternative. The model should remain advisory, with triage-nurse confirmation and a documented fallback.

The recommendation is conditional rather than permanent. It applies only to the
specified population, setting, intervention configuration, evidence cut-off, and
version. The following conditions should be incorporated into the approval,
contract, implementation, monitoring, or renewal record:

\begin{itemize}
 \item complete local calibration and silent-mode evaluation.
 \item demonstrate that rapid-assessment capacity can absorb expected alert volume.
 \item define subgroup and balancing measures, including delay among patients deprioritised by the system.
 \item cap total pilot expenditure and preserve the ability to stop without service disruption.
 \item review at 8 and 16 weeks using predefined proceed, redesign, restrict, or stop criteria.
\end{itemize}

\subsection{Decision Statement}
\begin{quote}
For the specified decision problem and comparator set, the institution should
\emph{[proceed / proceed with conditions / redesign / pilot / restrict / defer /
replace / retire]}. The decision applies to \emph{[population, setting, version,
and evidence period]}, is owned by \emph{[accountable authority]}, and will be
reviewed on \emph{[date]} or earlier if \emph{[trigger events]} occur. The
evidence, assumptions, unresolved uncertainty, mandatory safeguards, monitoring
thresholds, and exit arrangements are recorded in the accompanying dossier.
\end{quote}
